\lvli{PAT software design}
User Interface Design (UI-design) in Information Technology refers to the design of the user interface for software, websites or applications. It’s about programming the look of things, with a view to facilitating usability and to improve the user experience. Developing a user interface for complex systems simplifies the scientist’s work by giving him more time to focus on the experiment.

In this chapter, two analyses were carried out in order to build the PATSW part of this project. An analysis that focuses on the components available, and an analysis that focuses on what a user expects to use. Having these two analyses the construction of the single repositories becomes immediate, it is reduced to connecting the functionalities expected by users to the various components available.

\imgs{0.4}{pictures/ui}{Typica user interface for data analysis.}

\lvlii{The components analysis}
The aim of this project was to develop modular software that could control a PAT system. The PATSW had to be able to align the telescope with different light sources such as the light of a star or the laser of another telescope.
In addition, the PATSW had to be able to control all devices in the~\sect{Components}.
Starting from the PAT model~\fig{PAT} a simplified model has been developed~\fig{simplified-model} to highlight the various components to be implemented.
In this model we see that there are three main families to be implemented: sensors, controllers, and actuators. The arrows in this diagram indicate the information passages, in other words, actuator and sensor will have to communicate with the controller and vice versa.
It is therefore necessary to establish a common interface between the various sensors and actuators to connect them with the controllers. The study of these interfaces required multiple iterations to ensure maximum flexibility.
\img{simplified-model}{Simplified model for PAT system.}

\lvlii{The GUI analysis}
To get a complete idea of the goal of the system an other analysis was carry out. This analysis aims to establish the possible setups and what are the features that a user expects from the setup's GUI.
In this section will be reported the final interface that an user expects to use and its functionalities. This final interface it was created joining togheter the sub-modules interfaces.

\lvliii{Possible setups}
In this system there are three basic setups:
mirror-camera-pc, skywatcher-camera-pc, and mirror-PSD-KPA101.

\lvliv{Mirror-camera-pc}
In this setup the pc is connected with two devices: the first is
the KPA101 used to control the mirror; the second is the camera used as
sensor. Here two types of controllers are available: manual,
calibrated. Manual give the possibility to move the mirror using a joystick taking
in input times and voltages. Calibrated has two functionalities: the first is to use a joystick
similar to manual that it work taking in input a distance in pixels; the
second is to use and automatic controller that align the centroid to the
target that it is called One shot.

\lvliv{Skywatcher-camera-pc}
In this setup the pc is connected with two devices: the first is
the skywatcher mout; the second is the camera used as
sensor. The controllers available are the same of the mirror-camera-pc setup.

\lvliv{Mirror-PSD-KPA101}
In this setup the pc is connected only to the KPA101, the KPA101 is
then connected to the mirror and to the PSD. Here three types of controllers are available: manual, calibrated, PIDKPA101. Manual and calibrated are the same as the previous point. PIDKPA101 use the integrated PID-controller available in the KPA101.

\lvliii{Possible GUIs}

\img{gui}{Main GUI of the program.}

The GUI can be divided in two parts: the left side show the sensor
information, the right side show the controllers information.

In the sensor's side is composed by (form the top):

\begin{itemize}
      \item
            the frame rendering that contains three objects:
            \begin{itemize}

                  \item
                        the red dot represents the centroid of the beam;
                  \item
                        the red square represents the area where the controller calculates the
                        centroid;
                  \item
                        the green circle represents the target;
            \end{itemize}

      \item
            the set of parameters of the camera/PSD;
      \item
            the set of parameters to position the target;
      \item
            the plotter that show in real time the distances in pixels between the
            target and the centroid.
\end{itemize}

The controller's side there are three tabs: manual controller, calibrated
controller and PIDKPA101 controller (only for PSD).
The Manual controller and the PID KPA controller are simpler than calibrated.
In the manual controller voltages and times are the two parameters to be configured to optimize the joystick that will be use for the movement.
In the PIDKPA101 after the configuration of the PID parameters a checkbox \emph{enable PID}
is avaliable to start the controller. The \emph{auto open closed loop} checkbox is used to
disable the controller when the intensity of the beam go outside the
threshold range.
The calibrated controller is more complex so a separate section is reserved for it.

\lvliv{Calibrated controller}

\imgThree{0.3}{calibrated1}{Controller interface.}{0.3}{calibrated2}{Analyzer interface.}{0.3}{calibrated3}{Graph interface.}

The aim of this controller is to be able to use the \emph{One shot controller}, but before to be able to use it a calibration is required. The calibration procedure can be decomposed in two parts:
data acquisition, and data analysis.

\paragraph{Data acquisition}

In this phase several calibration points will be taken and stored into a file. The procedure
is the following:

\begin{itemize}
      \itemsep1pt\parskip0pt\parsep0pt
      \item
            in the \emph{Analyzer} tab set the voltages and the time interval;
      \item
            autogenerate the filename;
      \item
            start the point acquisition.
\end{itemize}

\paragraph{Data analysis}

In this phase the calibration points stored in a
file will be analyzed to retrieve the low/high voltages configuration. The procedure is the following:

\begin{itemize}
      \itemsep1pt\parskip0pt\parsep0pt
      \item
            in the \emph{Controller} tab select the file and press \emph{Upload
                  and Calibrated};
      \item
            go into \emph{Graphs} tab and check that the analysis is completed
            correctly;
      \item
            if everything is ok then go in the \emph{Controller} tab and set the
            configuration to the low voltages or high voltages pressing the
            corresponding button.
\end{itemize}

\paragraph{One shot controller}

After the selection of both low and high voltages, the One shot checkbox is ready to be pressed in the \emph{Controller} tab. To
swap between high and low voltages a time threshold is available in the same tab. If the movement is too long a high voltages will be selected to reduce the time, if not a low voltages will be selected to increase the precision.

\lvlii{The repositories structure}
\img{repositories}{This is a screenshot of all the repositories developed in the PAT group.}
Starting from the simplified model~\fig{simplified-model} has been chosen to create: a single repository for all controllers (Controller), a single repository for all actuators (Actuator), and two separate repositories for sensors (Camera, PSD).
As a communication interface between Controller and Actuator was made a Actuator class. As a communication interface between the Controller class and Camera/PSD it was chosen to create a Centroid class stored in a separate repository (Centroid).
The KPA101 repository was added to contain the KPA101 driver, while the Utils repository was created to contain a list of classes useful to all other repositories.
The Setup repository was created with the idea of taking pieces from the various repositories to compose different setups, in other words, it contains a list of possible PAT combinations.
The last two repositories that have been created are: Docs that contains quick documentation for programmers who want to contribute to the PAT, and PSD-bias that contains the software created to control a circuit that will be explained in the next chapter.

\lvlii{The repositories documentation}
This section lists all the documentations of the developed repositories.

      {\parindent0pt
            \lvliii{Utils}

This project contains a set of utilities for the PAT's repositories.
This REPO contains:

\begin{itemize}
  \tightlist
  \item
        Interval
  \item
        Point
  \item
        Range
  \item
        BoundedParameter
\end{itemize}

\textbf{Requirements}

\includegraphics[scale=0.7]{img/shilds/cpp.png}
\includegraphics[scale=0.7]{img/shilds/cmake.png}
\includegraphics[scale=0.7]{img/shilds/git.png}
\includegraphics[scale=0.7]{img/shilds/doxygen.png}
\includegraphics[scale=0.7]{img/shilds/sphinx.png}
\includegraphics[scale=0.7]{img/shilds/win.png}
\includegraphics[scale=0.7]{img/shilds/mac.png}
\includegraphics[scale=0.7]{img/shilds/linux.png}

\lvliv{Generality}

\lvlv{Import}

Import as an external library into your project by copy-paste the
following lines in your \texttt{config.json}.

\begin{lstlisting}[language=javascript, gobble=2]
  {
    "name"     : "PATUtils",
    "path"     : "gitlab.dei.unipd.it/PAT/Utils.git",
    "tag"      : "HEAD",
    "available": "YES",
    "getGui"   : "NO"
  }
\end{lstlisting}


\lvlv{Prerequisites}

The Prerequisites are the same descibed in the section "Setup your enviroment Prerequisites" \sect{Setup your environment}.

\lvliv{UML}
In figure \fig{utils} is illustrated the UML of the repository that shows the relationships between the classes that will be describe below.
\img{utils}{Utils UML.}

\lvliv{Usage}

\lvlv{Interval}

Interval give the possibility to generate a base-type structure that
has a \textbf{min} and \textbf{max} value.

First you need to include the header file and use the namespace.

\begin{lstlisting}[language=c++, gobble=2]
  #include <PAT/Utils/Interval.h>
  using namespace PAT::Utils;
\end{lstlisting}

Then you can define the variable ready to be used.

\begin{lstlisting}[language=c++, gobble=2]
  Interval<double> interval{0, 100};
  interval.min = 0;
  interval.max = 100;
\end{lstlisting}

\lvlv{Point}

Point give the possibility to generate a 2d-point that has a
\textbf{x} and \textbf{y} value.

First you need to include the header file and use the namespace.

\begin{lstlisting}[language=c++, gobble=2]
  #include <PAT/Utils/Point.h>
  using namespace PAT::Utils;
\end{lstlisting}

Then you can define the variable ready to be used.

\begin{lstlisting}[language=c++, gobble=2]
  Point<double> point{0, 100};
  point.x = 0;
  point.y = 100;  
\end{lstlisting}

It has some operations already defined suc as: sum, division, scalar
product.

\begin{lstlisting}[language=c++, gobble=2]
  Point<double> point2{0, 100};
  point += point2;
  point -= point2;
  point /= 3;
  point *= 3;  
\end{lstlisting}

\lvlv{Rectangle}

Rectangle give the possibility to generate a rectangular range defined
by two points(\textbf{Point}) that they are \textbf{start} and
\textbf{end}.

First you need to include the header file and use the namespace.

\begin{lstlisting}[language=c++, gobble=2]
  #include <PAT/Utils/Rectangle.h>
  using namespace PAT::Utils;
\end{lstlisting}

Then you can define the variable ready to be used.

\begin{lstlisting}[language=c++, gobble=2]
  Rectangle<double> rectangle{{0, 0}, {100, 100}};
  rectangle.start.x = 0;
  rectangle.start.y = 0;
  rectangle.end.x = 100;
  rectangle.end.y = 100;
  Point<double> point{50, 50};
  if(rectangle.isInside(point))
   rectangle.end = point;
\end{lstlisting}

\lvlv{Circle}

Circle give the possibility to generate a circular range defined by a
central point(Point) and a radius.

First you need to include the header file and use the namespace.

\begin{lstlisting}[language=c++, gobble=2]
  #include <PAT/Utils/Circle.h>
  using namespace PAT::Utils;
\end{lstlisting}

Then you can define the variable ready to be used.

\begin{lstlisting}[language=c++, gobble=2]
  Circle<double> circle{{0, 0}, 100};
  circle.center.x = 0;
  circle.center.y = 0;
  circle.radius = 100;
  Point<double> point{50, 50};
  if(circle.isInside(point))
    circle.radius = point.x;
  Rectangle<double> boundaries = circle.getBoundaries();
\end{lstlisting}

\lvlv{BoundedParameter}

BoundedParameter give the possibility to define a parameter that has:
\textbf{min} value, \textbf{max} value, \textbf{increment} value, and
\textbf{value} value.

First you need to include the header file and use the namespace.

\begin{lstlisting}[language=c++, gobble=2]
  #include <PAT/Utils/BoundedParameter.h>
  using namespace PAT::Utils;
\end{lstlisting}

Then you can define the variable ready to be used.

\begin{lstlisting}[language=c++, gobble=2]
  BoundedParameter<double> parameter{0, 100, 1, 50};
  parameter.min = 0;
  parameter.max = 100;
  parameter.increment = 1;
  parameter.value = 50;
\end{lstlisting}

\textbf{Warning}

The order of these parameters is imporant because is the same of the IDS
camera (min, max, increment).

            \lvliii{Centroid}

This project contains the centroid class. This class contains the
centroid information shared between the sensor and the controller in the
PAT-system.

\textbf{Requirements}

\includegraphics[scale=0.7]{img/shilds/cpp.png}
\includegraphics[scale=0.7]{img/shilds/cmake.png}
\includegraphics[scale=0.7]{img/shilds/git.png}
\includegraphics[scale=0.7]{img/shilds/doxygen.png}
\includegraphics[scale=0.7]{img/shilds/sphinx.png}
\includegraphics[scale=0.7]{img/shilds/win.png}
\includegraphics[scale=0.7]{img/shilds/mac.png}
\includegraphics[scale=0.7]{img/shilds/linux.png}

\lvliv{Generality}

\lvlv{Import}

Import as an external library into your project by copy-paste the
following lines in your \texttt{config.json}.

\begin{lstlisting}[language=javascript, gobble=2]
  {
    "name"     : "PATCentroid",
    "path"     : "gitlab.dei.unipd.it/PAT/Centroid.git",
    "tag"      : "HEAD",
    "available": "YES",
    "getGui"   : "NO"
  }
\end{lstlisting}

\lvlv{Prerequisites}

The Prerequisites are the same descibed in the section "Setup your enviroment Prerequisites" \sect{Setup your environment} with the addition of the following libraries that are auto fetched from the gitlab.dei.unipd.it:

\begin{itemize}
  \tightlist
  \item
        \href{https://gitlab.dei.unipd.it/PAT/Utils.git}{Utils} 1.0.0
\end{itemize}

\lvliv{UML}
In figure \fig{centroid} is illustrated the UML of the repository that shows the relationships between the classes that will be describe below.
\img{centroid}{Centroid UML.}

\lvliv{Usage}

Centroid give the possibility to store the point information with
\textbf{x} and \textbf{y} using a
\textbf{Point\textless double\textgreater{}} structure, and also the
\textbf{instensity} information.

It also contains a signal \textbf{newValue} that could be emitted
when a new value is ready.

First you need to include the header file and use the namespace.

\begin{lstlisting}[language=c++, gobble=2]
  #include <PAT/Centroid.h>
  using namespace PAT;
\end{lstlisting}

Then you can define the variable ready to be used.

\begin{lstlisting}[language=c++, gobble=2]
  Centroid centroid(10, 20, 100);
  centroid.x = 10;
  centroid.y = 20;
  centroid.instensity = 100;
\end{lstlisting}

If you need to emit the signal use this function.

\begin{lstlisting}[language=c++, gobble=2]
  emit centroid.newValue();
\end{lstlisting}

            \lvliii{KPA101}

This project create and interface between sensors, controllers,
actuators and all the functionalities of the KPA101.

\textbf{Requirements}

\includegraphics[scale=0.7]{img/shilds/cpp.png}
\includegraphics[scale=0.7]{img/shilds/cmake.png}
\includegraphics[scale=0.7]{img/shilds/git.png}
\includegraphics[scale=0.7]{img/shilds/doxygen.png}
\includegraphics[scale=0.7]{img/shilds/sphinx.png}
\includegraphics[scale=0.7]{img/shilds/win.png}

\lvlv{Generality}

\lvliv{Import}

Import as an external library into your project by copy-paste the
following lines in your \texttt{config.json}.

\begin{lstlisting}[language=javascript, gobble=2]
  {
    "name"     : "PATKPA101",
    "path"     : "gitlab.dei.unipd.it/PAT/KPA101.git",
    "tag"      : "HEAD",
    "available": "YES",
    "getGui"   : "NO"
  }
\end{lstlisting}

\lvliv{Prerequisites}

The following libraries are auto fetched from the gitlab.dei.unipd.it
host (ask the owner of this repo to become a member):

\begin{itemize}
  \tightlist
  \item
        \href{https://gitlab.dei.unipd.it/PAT/Utils.git}{Utils} 1.0.0
\end{itemize}

These other libraries need to be installed manually in your system:

\begin{itemize}
  \tightlist
  \item
        \href{https://www.qt.io/}{Qt} 5.14.2
\end{itemize}

The library documentation is generated through
\href{http://www.doxygen.nl/download.html}{Doxygen 1.8.13}. Additional
documentation in the \texttt{index} folder is generated through the
\href{https://www.anaconda.com/products/individual}{python3} package
\href{https://www.sphinx-doc.org/en/master/}{Sphinx} using the following
extensions (which you can install through pip3):

\begin{itemize}
  \tightlist
  \item
        \href{https://pypi.org/project/Sphinx/}{Sphinx 3.0.2}
  \item
        \href{https://sphinx-rtd-theme.readthedocs.io/en/stable/}{Sphinx read
          the doc theme} (to use the read the doc theme for html documentation)
  \item
        \href{https://pypi.org/project/breathe/}{Breathe} (to use the xml
        output of doxygen)
  \item
        \href{https://pypi.org/project/sphinx-markdown-builder/}{Sphinx-markdown-builder}
        (to generate the markdown version for gitlab wiki)
\end{itemize}

\texttt{pip3\ install\ Sphinx\ sphinx\_rtd\_theme\ breathe\ sphinx-markdown-builder}

\lvlv{Usage}

\lvliv{InputKPA}

InputKPA is a utility structure that contains the three voltagies that
KPA101 takes in input: \textbf{x}, \textbf{y}, \textbf{sum}.

First you need to include the header file and use the namespace.

\begin{lstlisting}[language=c++, gobble=2]
  #include <PAT/InputKPA.h>
  using namespace PAT;  
\end{lstlisting}

Then you can define the variable ready to be used.

\begin{lstlisting}[language=c++, gobble=2]
  InputKPA parameters;
  parameters.x = 1;
  parameters.y = 2;
  parameters.sum = 3;
\end{lstlisting}

\lvliv{KPA101}

KPA101 is a wrapper for the Kinesis libs that it enhance their
functionalities.

First you need to include the header file and use the namespace.

\begin{lstlisting}[language=c++, gobble=2]
  #include <PAT/KPA101.h>
  using namespace PAT;
\end{lstlisting}

In order to define a variable you can choose to pass the \textbf{kpa id}
or to let it automatically search it. If the search fails the program
throw an exception "not found". To search the id you can use
\textbf{find} and \textbf{findAll} functions.

\begin{lstlisting}[language=c++, gobble=2]
  KPA101 kpa101;          // auto
  KPA101 kpa101("kpaid"); // with id
  KPA101 kpa101(KPA101::find());       // with find
  KPA101 kpa101(KPA101::findAll()[0]); // with findAll
\end{lstlisting}

\textbf{Sensor}

Here are listed the functions useful for sensors.

If you need the get the current input state you can use this function.

\begin{lstlisting}[language=c++, gobble=2]
  InputKPA getInputVoltage() const;
\end{lstlisting}

\textbf{Actuator}

Here are listed the functions useful for actuators.

You can set the output state of the KPA using one of these three
functions. Output voltages are between 10 and 0.

\begin{itemize}
  \tightlist
  \item
        The function \textbf{setOutputVoltage} set the output voltages at the
        selected values.
  \item
        The function \textbf{setOutputImpulse} send a single impulse with
        length \textbf{\_time\_ms}.
  \item
        The function \textbf{setOutputCombinedImpulse} send two impulses at
        the same time with with their respective lengths. I one finish before
        the other its voltage go to zero.
\end{itemize}

\begin{lstlisting}[language=c++, gobble=2]
  double voltageX = 1;
  double voltageY = 1;
  kpa101.setOutputVoltage(voltageX, voltageY);

  double voltagePeakX = 1;
  double voltagePeakY = 1;
  double time_ms      = 1;
  kpa101.setOutputImpulse(voltagePeakX, voltagePeakY, time_ms);

  double voltagePeakX = 1;
  double timeX_ms     = 1;
  double voltagePeakY = 1;
  double timeY_ms     = 1;
  kpa101.setOutputCombinedImpulse(voltagePeakX, timeX_ms, voltagePeakY, timeY_ms);
\end{lstlisting}

\textbf{Controller}

Here are listed the functions useful to use the PID controller
integrated in the KPA.

You can set the PID parameters using these functions. PID values are
between 1 and 0.

\begin{lstlisting}[language=c++, gobble=2]
  auto parameters = kpa101.getPIDParameters();
  parameters.value.proportionalGain = 1;
  parameters.value.integralGain = 0.001;
  parameters.value.differentialGain = 0.01;
  kpa101.setPIDParameters(parameters);
\end{lstlisting}

You can enable/disable or check if enabled the PID controller using
these functions.

\begin{lstlisting}[language=c++, gobble=2]
  if(!kpa101.isPIDActive())
    kpa101.enablePID(true);
  else
    kpa101.enablePID(false);
\end{lstlisting}

You can enable/disable the auto open closed loop. If the intensity of
the light is outside the interval \textbf{autoOpenInterval} then the PID
controller is stopped. The intensity value is bounded between 10 and 0.

\begin{lstlisting}[language=c++, gobble=2]
  auto interval = kpa101.getAutoOpenInterval();
  interval.value.min = 0.1;
  interval.value.max = 1.3;
  kpa101.setAutoOpenInterval(interval);
  if(!kpa101.getAutoOpenCloseLoop())
    kpa101.setAutoOpenCloseLoop(true);
\end{lstlisting}

The output of the PID controller is multiplied by a feedback constant
that it range between -10 to +10 volts. Using this function you can set
it.

\begin{lstlisting}[language=c++, gobble=2]
  auto feedback = kpa101.getFeedBackGain();
  gain.value.x = 10;
  gain.value.y = 10;
  kpa101.setFeedBackGain(feedback);
\end{lstlisting}

\lvlv{UML}

\img{KPA101}{KPA101 UML.}
            \lvliii{Actuator}

This project contains the interface \textbf{Actuator} for all actuators. The
actual implementations are: STT254.

\textbf{Requirements}

\includegraphics[scale=0.7]{img/shilds/cpp.png}
\includegraphics[scale=0.7]{img/shilds/cmake.png}
\includegraphics[scale=0.7]{img/shilds/git.png}
\includegraphics[scale=0.7]{img/shilds/doxygen.png}
\includegraphics[scale=0.7]{img/shilds/sphinx.png}
\includegraphics[scale=0.7]{img/shilds/win.png}

\lvliv{Generality}

\lvlv{Import}

Import as an external library into your project by copy-paste the
following lines in your \texttt{config.json}.

\begin{lstlisting}[language=javascript, gobble=2]
  {
    "name"     : "PATActuator",
    "path"     : "gitlab.dei.unipd.it/PAT/Actuator.git",
    "tag"      : "HEAD",
    "available": "YES",
    "getGui"   : "NO"
  }
\end{lstlisting}

\lvlv{Prerequisites}

The following libraries are auto fetched from the gitlab.dei.unipd.it
host (ask the owner of this repo to become a member):

\begin{itemize}
  \tightlist
  \item
        \href{https://gitlab.dei.unipd.it/PAT/KPA101.git}{KPA101} 1.0.0
\end{itemize}

These other libraries need to be installed manually in your system:

\begin{itemize}
  \tightlist
  \item
        \href{https://www.qt.io/}{Qt} 5.14.2
\end{itemize}

The library documentation is generated through
\href{http://www.doxygen.nl/download.html}{Doxygen 1.8.13}. Additional
documentation in the \texttt{index} folder is generated through the
\href{https://www.anaconda.com/products/individual}{python3} package
\href{https://www.sphinx-doc.org/en/master/}{Sphinx} using the following
extensions (which you can install through pip3):

\begin{itemize}
  \tightlist
  \item
        \href{https://pypi.org/project/Sphinx/}{Sphinx 3.0.2}
  \item
        \href{https://sphinx-rtd-theme.readthedocs.io/en/stable/}{Sphinx read
          the doc theme} (to use the read the doc theme for html documentation)
  \item
        \href{https://pypi.org/project/breathe/}{Breathe} (to use the xml
        output of doxygen)
  \item
        \href{https://pypi.org/project/sphinx-markdown-builder/}{Sphinx-markdown-builder}
        (to generate the markdown version for gitlab wiki)
\end{itemize}

\texttt{pip3\ install\ Sphinx\ sphinx\_rtd\_theme\ breathe\ sphinx-markdown-builder}

\lvliv{Usage}

\lvlv{Actuator}

The \textbf{Actuator} class is a interface for all actuators. Here is the
include code.

\begin{lstlisting}[language=c++, gobble=2]
  #include <PAT/Actuator/Actuator.h>
\end{lstlisting}

To inherit the class you have to overwrite the two methods:
\textbf{setOutputImpulse}, \textbf{setOutputCombinedImpulse}. When the
impulse finish remember to emit a \textbf{movementCompleted} signal that
is already defined in \textbf{Actuator} class.

\begin{lstlisting}[language=c++, gobble=2]
  class YourClass : public Actuator {
    public:
      void setOutputImpulse(double _voltagePeakX, double _voltagePeakY, double _time_ms) override;
      void setOutputCombinedImpulse(double _voltagePeakX, double _timeX_ms, double _voltagePeakY, double _timeY_ms) override;
    };
\end{lstlisting}

\lvlv{STT254}

The \textbf{STT254} class is a Actuator implementation that works with
the \textbf{KPA101}. Here is the include code.

\begin{lstlisting}[language=c++, gobble=2]
  #include <PAT/KPA101.h>
  #include <PAT/Actuator/STT254.h>
  using namespace PAT::Actuator;
\end{lstlisting}

When you istantiate the \textbf{STT254} you must pass a \textbf{KPA101}
istance. The \textbf{STT254} is typically used as polymorphism with
\textbf{Actuator}.

\begin{lstlisting}[language=c++, gobble=2]
  std::shared_ptr<PAT::KPA101> kpa101         = std::make_shared<PAT::KPA101>();
  std::shared_ptr<PAT::Actuator::Actuator> actuator = std::make_shared<PAT::Actuator::STT254>(kpa101);
\end{lstlisting}

You can set the output state of the \textbf{Actuator} using one of these
two functions. Output voltages are between 10 and 0. When the functions
finish they emit a a \textbf{movementCompleted} signal.

\begin{itemize}
  \tightlist
  \item
        The function \textbf{setOutputImpulse} send a single impulse with
        length \textbf{\_time\_ms}.
  \item
        The function \textbf{setOutputCombinedImpulse} send two impulses at
        the same time with with their respective lengths. I one finish before
        the other its voltage go to zero.
\end{itemize}

\begin{lstlisting}[language=c++, gobble=2]
  double voltagePeakX = 1;
  double voltagePeakY = 1;
  double time_ms      = 1;
  actuator->setOutputImpulse(voltagePeakX, voltagePeakY, time_ms);
  
  double voltagePeakX = 1;
  double timeX_ms     = 1;
  double voltagePeakY = 1;
  double timeY_ms     = 1;
  actuator->setOutputCombinedImpulse(voltagePeakX, timeX_ms, voltagePeakY, timeY_ms);
\end{lstlisting}

\lvliv{UML}

\img{Mirror}{Actuator UML.}

            \lvliii{PSD}

This project contains the interface \textbf{PSD} for all PSD. The actual
implementations are: PDP90A.

\textbf{Requirements}

\includegraphics[scale=0.7]{img/shilds/cpp.png}
\includegraphics[scale=0.7]{img/shilds/cmake.png}
\includegraphics[scale=0.7]{img/shilds/git.png}
\includegraphics[scale=0.7]{img/shilds/doxygen.png}
\includegraphics[scale=0.7]{img/shilds/sphinx.png}
\includegraphics[scale=0.7]{img/shilds/win.png}

\lvliv{Generality}

\lvlv{Import}

Import as an external library into your project by copy-paste the
following lines in your \texttt{config.json}.

\begin{lstlisting}[language=javascript, gobble=2]
  {
    "name"     : "PATPSD",
    "path"     : "gitlab.dei.unipd.it/PAT/PSD.git",
    "tag"      : "HEAD",
    "available": "YES",
    "getGui"   : "YES"
  }
\end{lstlisting}

\lvlv{Prerequisites}

The following libraries are auto fetched from the gitlab.dei.unipd.it
host (ask the owner of this repo to become a member):

\begin{itemize}
  \tightlist
  \item
        \href{https://gitlab.dei.unipd.it/PAT/KPA101.git}{KPA101} 1.0.0
  \item
        \href{https://gitlab.dei.unipd.it/PAT/Centroid.git}{Centroid} 1.0.0
\end{itemize}

These other libraries need to be installed manually in your system:

\begin{itemize}
  \tightlist
  \item
        \href{https://www.qt.io/}{Qt} 5.14.2
\end{itemize}

The library documentation is generated through
\href{http://www.doxygen.nl/download.html}{Doxygen 1.8.13}. Additional
documentation in the \texttt{index} folder is generated through the
\href{https://www.anaconda.com/products/individual}{python3} package
\href{https://www.sphinx-doc.org/en/master/}{Sphinx} using the following
extensions (which you can install through pip3):

\begin{itemize}
  \tightlist
  \item
        \href{https://pypi.org/project/Sphinx/}{Sphinx 3.0.2}
  \item
        \href{https://sphinx-rtd-theme.readthedocs.io/en/stable/}{Sphinx read
          the doc theme} (to use the read the doc theme for html documentation)
  \item
        \href{https://pypi.org/project/breathe/}{Breathe} (to use the xml
        output of doxygen)
  \item
        \href{https://pypi.org/project/sphinx-markdown-builder/}{Sphinx-markdown-builder}
        (to generate the markdown version for gitlab wiki)
\end{itemize}

\texttt{pip3\ install\ Sphinx\ sphinx\_rtd\_theme\ breathe\ sphinx-markdown-builder}

\lvliv{Usage}

\lvlv{PSD}

The \textbf{PSD} class is a interface for all PSD. Here is the include
code.

\begin{lstlisting}[language=c++, gobble=2]
  #include <PAT/PSD/PSD.h>
\end{lstlisting}

To inherit the class you have to overwrite the three methods:
\textbf{start}, \textbf{stop}, \textbf{getIntensity}.

\begin{lstlisting}[language=c++, gobble=2]
  class YourClass : public PSD {
    public:
      void start() override;
      void stop() override;
      double getIntensity() const override;
    };    
\end{lstlisting}

\lvlv{PDP90A}

The \textbf{PDP90A} class is a PSD implementation that it works with the
\textbf{KPA101}. Here is the include code.

\begin{lstlisting}[language=c++, gobble=2]
  #include <PAT/KPA101.h>
  #include <PAT/PSD/PDP90A.h>
  using namespace PAT::PSD;  
\end{lstlisting}

When you istantiate the \textbf{PDP90A} you must pass a \textbf{KPA101}
istance. The \textbf{PDP90A} is typically used as polymorphism with
\textbf{PSD}.

\begin{lstlisting}[language=c++, gobble=2]
  std::shared_ptr<PAT::KPA101> kpa101   = std::make_shared<PAT::KPA101>();
  std::shared_ptr<PAT::PSD::PSD> psd = std::make_shared<PAT::PSD::PDP90A>(kpa101);
\end{lstlisting}

You can enable/disable or check if enabled the PDP90A's acquisition
using these functions.

\begin{lstlisting}[language=c++, gobble=2]
  if(!psd->isActive())
    psd->start();
 else
    psd->stop();
\end{lstlisting}

You can set the frame rate of the acquisition using these functions. At
every frame the PSD calculates the centroid and then it emit a signal
\textbf{newFrameSignal}.

\begin{lstlisting}[language=c++, gobble=2]
  auto framerate  = psd->getFrameRate();
  framerate.value = 30;
  psd->setFrameRate(framerate);
\end{lstlisting}

You can get the other parameters using these functions.

\begin{lstlisting}[language=c++, gobble=2]
  PAT::Utils::Point<double> sensorSize = psd->getSensorSize();
  std::shared_ptr<PAT::Utils::BoundedParameter<PAT::Utils::Circle<int>>> target = psd->getTarget();
  std::shared_ptr<PAT::Centroid> centroid = psd->getCentroid();
  PAT::Utils::Point<double> sensorSize = psd->getSensorSize();
  double intensity = psd->getIntensity();
\end{lstlisting}

\lvlv{PSDWiget}

The \textbf{PSDWiget} class is a widget fot the PSD's implementations.
Here I present the PDP90A example. The include code is the following.

\begin{lstlisting}[language=c++, gobble=2]
  #include <QApplication>
  #include <PAT/PSD/gui/PSDWidget.h>
  #include <PAT/PSD/PDP90A.h>
  using namespace PAT::PSD;
\end{lstlisting}

This is a simple \textbf{main} example.

\begin{lstlisting}[language=c++, gobble=2]
  int main(int argc, char *argv[]) {
    QApplication a(argc, argv);
    PSDWidget cameraWidget;
    std::shared_ptr<PAT::KPA101> kpa101 = std::make_shared<PAT::KPA101>();
    std::shared_ptr<PSD> camera         = std::make_shared<PDP90A>(kpa101);
    cameraWidget.setPSD(camera);
    cameraWidget.show();
    camera->start();
    return a.exec();
 }
\end{lstlisting}

If you want instead integrate the widget into another GUI (layout) you
can use this code.

\begin{lstlisting}[language=c++, gobble=2]
  QWidget *window                     = new QWidget();
  QVBoxLayout *mainLayout             = new QVBoxLayout();
  PSDWidget *psdWidget                = new PSDWidget();
  std::shared_ptr<PAT::KPA101> kpa101 = std::make_shared<PAT::KPA101>();
  std::shared_ptr<PSD> camera         = std::make_shared<PDP90A>(kpa101);
  
  psdWidget->setPSD(psd);
  mainLayout->addWidget((QWidget *)psdWidget);
  window->setLayout(mainLayout);
  window->show();
\end{lstlisting}

\lvliv{UML}

\img{PSD}{PSD UML.}

            \lvliii{Camera}

This project contains the interface \textbf{Camera} for all cameras. The
actual implementations are: IDS.

\textbf{Requirements}

\includegraphics[scale=0.7]{img/shilds/cpp.png}
\includegraphics[scale=0.7]{img/shilds/cmake.png}
\includegraphics[scale=0.7]{img/shilds/git.png}
\includegraphics[scale=0.7]{img/shilds/doxygen.png}
\includegraphics[scale=0.7]{img/shilds/sphinx.png}
\includegraphics[scale=0.7]{img/shilds/win.png}

\lvliv{Generality}

\lvlv{Import}

Import as an external library into your project by copy-paste the
following lines in your \texttt{config.json}.

\begin{lstlisting}[language=javascript, gobble=2, caption=PATCamera module inclusion.]
  {
    "name"     : "PATCamera",
    "path"     : "gitlab.dei.unipd.it/PAT/Camera.git",
    "tag"      : "HEAD",
    "available": "YES",
    "getGui"   : "YES"
  }
\end{lstlisting}

\lvlv{Prerequisites}

The following libraries are auto fetched from the gitlab.dei.unipd.it
host (ask the owner of this repo to become a member):

\begin{itemize}
  \tightlist
  \item
        \href{https://gitlab.dei.unipd.it/PAT/Centroid.git}{Centroid} 1.0.0
\end{itemize}

These other libraries need to be installed manually in your system:

\begin{itemize}
  \tightlist
  \item
        \href{https://www.qt.io/}{Qt} 5.14.2
\end{itemize}

The library documentation is generated through
\href{http://www.doxygen.nl/download.html}{Doxygen 1.8.13}. Additional
documentation in the \texttt{index} folder is generated through the
\href{https://www.anaconda.com/products/individual}{python3} package
\href{https://www.sphinx-doc.org/en/master/}{Sphinx} using the following
extensions (which you can install through pip3):

\begin{itemize}
  \tightlist
  \item
        \href{https://pypi.org/project/Sphinx/}{Sphinx 3.0.2}
  \item
        \href{https://sphinx-rtd-theme.readthedocs.io/en/stable/}{Sphinx read
          the doc theme} (to use the read the doc theme for html documentation)
  \item
        \href{https://pypi.org/project/breathe/}{Breathe} (to use the xml
        output of doxygen)
  \item
        \href{https://pypi.org/project/sphinx-markdown-builder/}{Sphinx-markdown-builder}
        (to generate the markdown version for gitlab wiki)
\end{itemize}

\texttt{pip3\ install\ Sphinx\ sphinx\_rtd\_theme\ breathe\ sphinx-markdown-builder}

\lvliv{Usage}

\lvlv{Frame}

The \textbf{Frame} class manages all the processing of the images
captured by the camera.

\begin{lstlisting}[language=c++, gobble=2, caption=PAT/Camera/Frame inclusion.]
  #include <PAT/Camera/Frame.h>
\end{lstlisting}

This is an example of initialization.

\begin{lstlisting}[language=c++, gobble=2, caption=Frame initialization example.]
  int width = 100;
  int height = 100;
  unsigned char imgBuffer[width * height] = {0};
  std::shared_ptr<Frame> frame = std::make_shared<Frame>(imgBuffer, width, height, QImage::Format_Grayscale8);
\end{lstlisting}

You can flip the image in the two axes usign these functions.

\begin{lstlisting}[language=c++, gobble=2]
  frame->flipX();
  frame->flipY();
\end{lstlisting}

You can calculate the centroid, the spotDiameter and the integral (spot
intensity) using this function. If it does not find the centroid it
returns -1 as invalid value.

\begin{lstlisting}[language=c++, gobble=2]
  frame->calculateCentroidAndIntegral();
\end{lstlisting}

You can also control the centroid calculation setting the parameters
\textbf{threshold} and \textbf{rangeRadius}. -
\textbf{setThreshold/getThreshold} set/get the minimum value for pixels
to be counted in the centroid average. -
\textbf{setRangeRadius/getRangeRadius} set/get the range radius in pixel
where to search the centroid. The value 0 means infinite range. -
\textbf{getRange} get the four cordinate of the range where to search
the centroid.

\begin{lstlisting}[language=c++, gobble=2]
  auto threshold = frame->getThreshold();
  threshold.value = 30;
  frame->setThreshold(threshold);
  
  auto rangeRadius = frame->getRangeRadius();
  rangeRadius.value = 200;
  frame->setRangeRadius(rangeRadius);
  
  PAT::Utils::Range<int> range = frame->getRange();
\end{lstlisting}

You can get the other parameters using these functions.

\begin{lstlisting}[language=c++, gobble=2]
  std::shared_ptr<PAT::Centroid> centroid = frame->getCentroid();
  std::shared_ptr<PAT::Utils::BoundedParameter<PAT::Utils::Circle<int>>> target = frame->getTarget();
  unsigned char* imgBuffer = frame->getImageRaw();
  QImage img = frame->getImage();
  int width = frame->getWidth();
  int height = frame->getHeight();
\end{lstlisting}

\lvlv{Camera}

The \textbf{Camera} class is a interface for all Camera. Here is the
include code.

\begin{lstlisting}[language=c++, gobble=2]
  #include <PAT/Camera/Camera.h>
\end{lstlisting}

To inherit the class you have to overwrite the four methods:
\textbf{start}, \textbf{stop}, \textbf{setFrameRate},
\textbf{setExposureTime}.

\begin{lstlisting}[language=c++, gobble=2]
  class YourClass : public Camera {
    public:
      void start() override;
      void stop() override;
      void setFrameRate(const PAT::Utils::BoundedParameter<double>& _frameRate) override;
      void setExposureTime(const PAT::Utils::BoundedParameter<double>& _exposureTimeValue) override;
    };
\end{lstlisting}

\lvlv{IDS}

The \textbf{IDS} class is a Camera implementation. Here is the include
code.

\begin{lstlisting}[language=c++, gobble=2]
  #include <PAT/Camera/IDS.h>
  using namespace PAT::Camera;
\end{lstlisting}

The \textbf{IDS} is typically used as polymorphism with \textbf{Camera}.
In order to define a variable you can choose to pass the \textbf{IDS id}
or to let it automatically search it. If the search fails the program
throw an exception "not found". To search the id you can use
\textbf{find} and \textbf{findAll} functions.

\begin{lstlisting}[language=c++, gobble=2]
  std::shared_ptr<PAT::Camera::Camera> camera = std::make_shared<PAT::Camera::IDS>();
  std::shared_ptr<PAT::Camera::Camera> camera = std::make_shared<PAT::Camera::IDS>(IDS::find());
  std::shared_ptr<PAT::Camera::Camera> camera = std::make_shared<PAT::Camera::IDS>(IDS::findAll()->uci[0].dwCameraID);
\end{lstlisting}

You can enable/disable or check if enabled the IDS's acquisition using
these functions.

\begin{lstlisting}[language=c++, gobble=2]
  if(!camera->isActive())
    camera->start();
 else
    camera->stop();
\end{lstlisting}

You can set the frame rate of the acquisition using these functions. At
every frame the Camera calculates the centroid and then it emit a signal
\textbf{newFrameSignal}.

\begin{lstlisting}[language=c++, gobble=2]
  auto framerate  = camera->getFrameRate();
  framerate.value = 30;
  camera->setFrameRate(framerate);
\end{lstlisting}

You can set the exposure time of the acquisition using these functions.
The exposure time limit depends on the frame rate (1 / framerate).

\begin{lstlisting}[language=c++, gobble=2]
  auto exposureTime  = camera->getExposureTime();
  exposureTime.value = 1 / 30;
  camera->setExposureTime(exposureTime);
\end{lstlisting}

You can enable/disable or check if enabled the centroid calculation
using these functions.

\begin{lstlisting}[language=c++, gobble=2]
  if(!camera->getCentroidCalculation())
    camera->setCentroidCalculation(true);
 else
    camera->setCentroidCalculation(false);
\end{lstlisting}

You can alse obtain the frame object that it automatically update at
every frame.

\begin{lstlisting}[language=c++, gobble=2]
  std::shared_ptr<Frame> frame = camera->getFrame();
\end{lstlisting}

\lvlv{CameraWiget}

The \textbf{CameraWiget} class is a widget fot the Camera's
implementations. Here I present the IDS example. The include code is the
following.

\begin{lstlisting}[language=c++, gobble=2]
  #include <QApplication>
  #include <PAT/Camera/gui/CameraWidget.h>
  #include <PAT/Camera/IDS.h>
  using namespace PAT::Camera;
\end{lstlisting}

This is a simple \textbf{main} example.

\begin{lstlisting}[language=c++, gobble=2]
  int main(int argc, char *argv[]) {
    QApplication a(argc, argv);
    CameraWidget cameraWidget;
    std::shared_ptr<Camera> camera = std::make_shared<IDS>();
    cameraWidget.setCamera(camera);
    cameraWidget.show();
    camera->start();
   return a.exec();
  }
\end{lstlisting}

If you want instead integrate the widget into another GUI (layout) you
can use this code.

\begin{lstlisting}[language=c++, gobble=2]
  QWidget *window                = new QWidget();
  QVBoxLayout *mainLayout        = new QVBoxLayout();
  CameraWidget *cameraWidget     = new CameraWidget();
  std::shared_ptr<Camera> camera = std::make_shared<IDS>();
  
  cameraWidget->setCamera(camera);
  mainLayout->addWidget((QWidget *)cameraWidget);
  window->setLayout(mainLayout);
  window->show();
\end{lstlisting}

\lvliv{UML}

\img{Camera}{Camera UML.}

            \lvliii{Controller}

This project contains different controllers for the PAT system. The
controllers implemented at moment are: Manual, Calibrated, PIDKPA101.

\textbf{Requirements}

\includegraphics[scale=0.7]{img/shilds/cpp.png}
\includegraphics[scale=0.7]{img/shilds/cmake.png}
\includegraphics[scale=0.7]{img/shilds/git.png}
\includegraphics[scale=0.7]{img/shilds/doxygen.png}
\includegraphics[scale=0.7]{img/shilds/sphinx.png}
\includegraphics[scale=0.7]{img/shilds/win.png}

\lvliv{Generality}

\lvlv{Import}

Import as an external library into your project by copy-paste the
following lines in your \texttt{config.json}.

\begin{lstlisting}[language=javascript, gobble=2]
    {
      "name"     : "PATController",
      "path"     : "gitlab.dei.unipd.it/PAT/Controller.git",
      "tag"      : "HEAD",
      "available": "YES",
      "getGui"   : "YES"
    }
\end{lstlisting}

\lvlv{Prerequisites}

The following libraries are auto fetched from the gitlab.dei.unipd.it
host (ask the owner of this repo to become a member):

\begin{itemize}
  \tightlist
  \item
        \href{https://gitlab.dei.unipd.it/PAT/Mirror}{Mirror} 1.0.0
  \item
        \href{https://gitlab.dei.unipd.it/PAT/PSD}{PSD} 1.0.0
  \item
        \href{https://gitlab.dei.unipd.it/PAT/Camera}{Camera} 1.0.0
\end{itemize}

These other libraries need to be installed manually in your system:

\begin{itemize}
  \tightlist
  \item
        \href{https://www.qt.io/}{Qt} 5.14.2
\end{itemize}

The library documentation is generated through
\href{http://www.doxygen.nl/download.html}{Doxygen 1.8.13}. Additional
documentation in the \texttt{index} folder is generated through the
\href{https://www.anaconda.com/products/individual}{python3} package
\href{https://www.sphinx-doc.org/en/master/}{Sphinx} using the following
extensions (which you can install through pip3):

\begin{itemize}
  \tightlist
  \item
        \href{https://pypi.org/project/Sphinx/}{Sphinx 3.0.2}
  \item
        \href{https://sphinx-rtd-theme.readthedocs.io/en/stable/}{Sphinx read
          the doc theme} (to use the read the doc theme for html documentation)
  \item
        \href{https://pypi.org/project/breathe/}{Breathe} (to use the xml
        output of doxygen)
  \item
        \href{https://pypi.org/project/sphinx-markdown-builder/}{Sphinx-markdown-builder}
        (to generate the markdown version for gitlab wiki)
\end{itemize}

\texttt{pip3\ install\ Sphinx\ sphinx\_rtd\_theme\ breathe\ sphinx-markdown-builder}

\lvliv{Usage}

\lvlv{PIDKPA101}

This class offers a set of functions to comunicate with the PID
controller integrated in the KPA101.

First you need to include the header file and use the namespace.

\begin{lstlisting}[language=c++, gobble=2]
  #include <PAT/Controller/PIDKPA101.h>
  using namespace PAT::Controller;
\end{lstlisting}

When you istantiate the \textbf{PIDKPA101} you must pass a
\textbf{KPA101} istance.

\begin{lstlisting}[language=c++, gobble=2]
  std::shared_ptr<PAT::KPA101> kpa101   = std::make_shared<PAT::KPA101>();
  std::shared_ptr<PAT::Controller::PIDKPA101> pidkpa101 = std::make_shared<PAT::Controller::PIDKPA101>(kpa101);
\end{lstlisting}


You can enable/disable or check if enabled the PIDKPA101 controller
using these functions.

\begin{lstlisting}[language=c++, gobble=2]
  if(!pidkpa101->isPIDActive())
    pidkpa101->enablePID(true);
 else
    pidkpa101->enablePID(false); 
\end{lstlisting}


You can set the PID parameters using these functions. PID values are
between 1 and 0.

\begin{lstlisting}[language=c++, gobble=2]
  auto parameters = pidkpa101->getPIDParameters();
  parameters.value.proportionalGain = 1;
  parameters.value.integralGain = 0.001;
  parameters.value.differentialGain = 0.01;
  pidkpa101->setPIDParameters(parameters);
\end{lstlisting}


You can enable/disable the auto open closed loop. If the intensity of
the light is outside the interval \textbf{autoOpenInterval} then the PID
controller is stopped. The intensity value is bounded between 10 and 0.

\begin{lstlisting}[language=c++, gobble=2]
  auto interval = pidkpa101->getAutoOpenInterval();
  interval.value.min = 0.1;
  interval.value.max = 1.3;
  pidkpa101->setAutoOpenInterval(interval);
  if(!pidkpa101->getAutoOpenCloseLoop())
   pidkpa101->setAutoOpenCloseLoop(true);
\end{lstlisting}

The output of the PID controller is multiplied by a feedback constant
that it range between -10 to +10 volts. Using this function you can set
it.

\begin{lstlisting}[language=c++, gobble=2]
  auto feedback = pidkpa101->getFeedBackGain();
  gain.value.x = 10;
  gain.value.y = 10;
  pidkpa101->setFeedBackGain(feedback);
\end{lstlisting}


\lvlv{Manual}

Manual is a controller that it allow to move a mirror in the four
directions. It require a \textbf{Mirror} and a \textbf{Centroid}. This
last one is from a \textbf{Frame} inside a \textbf{Camera} or from a
\textbf{PSD}. For semplicity here I provide an example using a
\textbf{IDS} camera and a \textbf{STT254} mirror.

First you need to include the header file and use the namespace.

\begin{lstlisting}[language=c++, gobble=2]
  #include <PAT/Controller/Manual.h>
  #include <PAT/Mirror/STT254.h>
  #include <PAT/Camera/IDS.h>
  using namespace PAT::Controller;
  using namespace PAT::Camera;
  using namespace PAT::Mirror;
\end{lstlisting}


When you istantiate the \textbf{Manual} you must pass a \textbf{Mirror}
and a \textbf{Centroid} istances.

\begin{lstlisting}[language=c++, gobble=2]
  std::shared_ptr<Camera> camera      = std::make_shared<IDS>();
  std::shared_ptr<PAT::KPA101> kpa101 = std::make_shared<PAT::KPA101>();
  std::shared_ptr<Mirror> mirror      = std::make_shared<STT254>(kpa101);
  std::shared_ptr<Manual> manual      = std::make_shared<Manual>(mirror, camera->getFrame()->getCentroid());
\end{lstlisting}


To move in one of the four directions you can use this code.

\begin{lstlisting}[language=c++, gobble=2]
  manual->move(Direction::up);
  manual->move(Direction::down);
  manual->move(Direction::left);
  manual->move(Direction::right);
\end{lstlisting}


To choose how much move you need to set times (milliseconds) and
voltages (volts) in the four directions. You can also control the
current values.

\begin{lstlisting}[language=c++, gobble=2]
  if(manual->getManualVoltage(Direction::up) > 100)
    manual->setManualVoltage(1, Direction::up);
  if(manual->getManualTime(Direction::up) > 100)
    manual->setManualTime(1, Direction::down);
\end{lstlisting}


You can also invert or check the movement directions in the two axes.
You can also swap the two axes.

\begin{lstlisting}[language=c++, gobble=2]
  if(!manual->getFlippedX())
    manual->setFlippedX(true);
  else
    manual->setFlippedX(false);
  
  if(!manual->getFlippedY())
    manual->setFlippedY(true);
  else
    manual->setFlippedY(false);
  
  if(!manual->getSwappedXY())
    manual->setSwappedXY(true);
  else
    manual->setSwappedXY(false);
\end{lstlisting}


\lvlv{ManualWidget}

The \textbf{ManualWidget} class is a widget for the Manual controller.
Here I present the an example consistent with the Manual one. The
include code is the following.

\begin{lstlisting}[language=c++, gobble=2]
  #include <QApplication>
  #include <PAT/Controller/gui/ManualWidget.h>
  #include <PAT/Mirror/STT254.h>
  #include <memory>
  #include <PAT/Camera/IDS.h>
  using namespace PAT::Controller;
  using namespace PAT::Camera;
  using namespace PAT::Mirror;
  
\end{lstlisting}


This is a simple \textbf{main} example.

\begin{lstlisting}[language=c++, gobble=2]
  int main(int argc, char *argv[]) {
    QApplication a(argc, argv);
    ManualWidget manualWidget;
    std::shared_ptr<Camera> camera      = std::make_shared<IDS>();
    std::shared_ptr<PAT::KPA101> kpa101 = std::make_shared<PAT::KPA101>();
    std::shared_ptr<Mirror> mirror      = std::make_shared<STT254>(kpa101);
    std::shared_ptr<Manual> manual      = std::make_shared<Manual>(mirror, camera->getFrame()->getCentroid());
    manualWidget.setController(manual);
    manualWidget.show();
    camera->start();
    return a.exec();
  }
\end{lstlisting}


If you want instead integrate the widget into another GUI (layout) you
can use this code.

\begin{lstlisting}[language=c++, gobble=2]
  QWidget *window                     = new QWidget();
  QVBoxLayout *mainLayout             = new QVBoxLayout();
  ManualWidget *manualWidget          = new ManualWidget();
  
  std::shared_ptr<Camera> camera      = std::make_shared<IDS>();
  std::shared_ptr<PAT::KPA101> kpa101 = std::make_shared<PAT::KPA101>();
  std::shared_ptr<Mirror> mirror      = std::make_shared<STT254>(kpa101);
  std::shared_ptr<Manual> manual      = std::make_shared<Manual>(mirror, camera->getFrame()->getCentroid());
  
  manualWidget->setPSD(manual);
  mainLayout->addWidget((QWidget *)manualWidget);
  window->setLayout(mainLayout);
  window->show();
\end{lstlisting}


There are also these three signal functions that they are used to
propagate the click of their respective buttons in other controllers.

\begin{lstlisting}[language=c++, gobble=2]
  void chkFlipXchanged(int arg1);
  void chkFlipYchanged(int arg1);
  void chkSwapXYchanged(int arg1);
\end{lstlisting}


\lvlv{CalibrationConfig}

This structure contains the main parameters for the calibration. The
parameters are: the four voltages, one per direction; the time durations
of the movements described by min-max-increment; the inversion or swap
of the axes. The include code is the following.

\begin{lstlisting}[language=c++, gobble=2]
  #include <PAT/Controller/CalibrationConfig.h>
\end{lstlisting}


You can also store/unstore from a file these parameters.

\begin{lstlisting}[language=c++, gobble=2]
  CalibrationConfig parameters;
  QString filename = "test.txt";
  parameters.write(filename);
  parameters.read(filename);
\end{lstlisting}


\lvlv{CalibrationAnalysis}

This is a support structure to take the data from a file and fit them
using a linear fit. The include code is the following.

\begin{lstlisting}[language=c++, gobble=2]
  #include <PAT/Controller/CalibrationAnalysis.h>
\end{lstlisting}

This structure can analyze file data generated by \textbf{Calibrated}
controller.

\begin{lstlisting}[language=c++, gobble=2]
  CalibrationAnalysis analysis;
  QString filename = "test.txt";
  analysis.read(filename);
\end{lstlisting}

\lvlv{Calibrated}

Calibrated is a controller that it allow to move a mirror in the four
directions. It is able to calibrate itself using linear fit. It also
contain a One shot controller that in one frame it try to align the
centroid to the target point. It require a \textbf{Mirror}, a
\textbf{Centroid} and a target point. The \textbf{Centroid} and the
calibrated point are from a \textbf{Frame} inside a \textbf{Camera} or
from a \textbf{PSD}. For semplicity here I provide an example using a
\textbf{IDS} camera and a \textbf{STT254} mirror.

First you need to include the header file and use the namespace.

\begin{lstlisting}[language=c++, gobble=2]
  #include <PAT/Controller/Calibrated.h>
  #include <PAT/Mirror/STT254.h>
  #include <PAT/Camera/IDS.h>
  using namespace PAT::Controller;
  using namespace PAT::Camera;
  using namespace PAT::Mirror;
\end{lstlisting}

When you istantiate the \textbf{Calibrated} you must pass a
\textbf{Mirror} and a \textbf{Centroid} istances.

\begin{lstlisting}[language=c++, gobble=2]
  std::shared_ptr<Camera> camera         = std::make_shared<IDS>();
  std::shared_ptr<PAT::KPA101> kpa101    = std::make_shared<PAT::KPA101>();
  std::shared_ptr<Mirror> mirror         = std::make_shared<STT254>(kpa101);
  std::shared_ptr<Calibrated> calibrated = std::make_shared<Calibrated>(mirror, camera->getFrame()->getCentroid(), camera->getFrame()->getTarget());
\end{lstlisting}

To calibrate the controller you need first to collect the data. To do so
you need to choose a configuration and a file name. You can use this
code as example of command you require.

\begin{lstlisting}[language=c++, gobble=2]
  auto upVoltage     = calibrated->getInputVoltage(Direction::up);
  auto downVoltage   = calibrated->getInputVoltage(Direction::down);
  auto leftVoltage   = calibrated->getInputVoltage(Direction::left);
  auto rightVoltage  = calibrated->getInputVoltage(Direction::right);
  upVoltage.value    = 2;
  downVoltage.value  = 2;
  leftVoltage.value  = 2;
  rightVoltage.value = 2;
  calibrated->setInputVoltage(upVoltage, Direction::up);
  calibrated->setInputVoltage(downVoltage, Direction::down);
  calibrated->setInputVoltage(leftVoltage, Direction::left);
  calibrated->setInputVoltage(rightVoltage, Direction::right);
  
  auto timeMin  = calibrated->getTimeIntervalMin();
  auto timeInc  = calibrated->getTimeIntervalInc();
  auto timeMax  = calibrated->getTimeIntervalMax();
  timeMin.value = 1;
  timeInc.value = 0.5;
  timeMax.value = 5;
  calibrated->setTimeIntervalMin(timeMin);
  calibrated->setTimeIntervalInc(timeInc);
  calibrated->setTimeIntervalMax(timeMax);
  
  auto nPoints = calibrated->getNPoints();
  nPoints.value = 10;
  calibrated->setNPoints(nPoints);
  
  calibrated->setFlippedX(true);
  calibrated->setFlippedY(false);
  calibrated->setSwappedXY(false);
  
  calibrated->autoFileName();
\end{lstlisting}

After selected the config you can start/stop to collect the data.

\begin{lstlisting}[language=c++, gobble=2]
  calibrated->startTakeCalibrationPoints();
  calibrated->stopTakeCalibrationPoints();
\end{lstlisting}

After collected the data you can analyze it and store the result in
cache.

\begin{lstlisting}[language=c++, gobble=2]
  calibrated->analyzeCalibrationPoints();
\end{lstlisting}

To conclude the process, when you are happy with the cached data you can
use it as configuration to the calibrated movement. Remember to select
the time threshold in milliseconds. If the calibrated movement time is
higher than threshold the controller use an high voltage else a low
voltage.

\begin{lstlisting}[language=c++, gobble=2]
  calibrated->useCurrentCalibration(VoltageType::low);
  calibrated->useCurrentCalibration(VoltageType::high);
  auto threshold = calibrated->getThreshold();
  threshold.value = 2;
  calibrated->setThreshold(threshold);
\end{lstlisting}

Completed this process you can use the directional movement or the One
Shot controller.

\begin{lstlisting}[language=c++, gobble=2]
  auto deltaPx = calibrated->getDeltaPixel();
  deltaPx.value = 50;
  calibrated->setDeltaPixel(deltaPx);
  calibrated->moveAuto(Direction::up);
  
  calibrated->startOneShotController();
  calibrated->stopOneShotController();
\end{lstlisting}

\lvlv{CalibratedWidget}

The \textbf{CalibratedWidget} class is a widget for the Calibrated
controller. Here I present the an example consistent with the Calibrated
one. The include code is the following.

\begin{lstlisting}[language=c++, gobble=2]
  #include <QApplication>
  #include <PAT/Controller/gui/CalibratedWidget.h>
  #include <PAT/Mirror/STT254.h>
  #include <memory>
  #include <PAT/Camera/IDS.h>
  using namespace PAT::Controller;
  using namespace PAT::Camera;
  using namespace PAT::Mirror;
\end{lstlisting}

This is a simple \textbf{main} example.

\begin{lstlisting}[language=c++, gobble=2]
  int main(int argc, char *argv[]) {
    QApplication a(argc, argv);
    CalibratedWidget calibratedWidget;
    std::shared_ptr<Camera> camera      = std::make_shared<IDS>();
    std::shared_ptr<PAT::KPA101> kpa101 = std::make_shared<PAT::KPA101>();
    std::shared_ptr<Mirror> mirror      = std::make_shared<STT254>(kpa101);
    std::shared_ptr<Calibrated> calibrated  = std::make_shared<Calibrated>(mirror, camera->getFrame()->getCentroid(), camera->getFrame()->getTarget());
    calibratedWidget.setController(calibrated);
    calibratedWidget.show();
    camera->start();
    return a.exec();
  }
\end{lstlisting}

If you want instead integrate the widget into another GUI (layout) you
can use this code.

\begin{lstlisting}[language=c++, gobble=2]
  QWidget *window                     = new QWidget();
  QVBoxLayout *mainLayout             = new QVBoxLayout();
  CalibratedWidget *calibratedWidget  = new CalibratedWidget();
  
  std::shared_ptr<Camera> camera         = std::make_shared<IDS>();
  std::shared_ptr<PAT::KPA101> kpa101    = std::make_shared<PAT::KPA101>();
  std::shared_ptr<Mirror> mirror         = std::make_shared<STT254>(kpa101);
  std::shared_ptr<Calibrated> calibrated = std::make_shared<Calibrated>(mirror, camera->getFrame()->getCentroid());
  
  calibratedWidget->setPSD(calibrated);
  mainLayout->addWidget((QWidget *)calibratedWidget);
  window->setLayout(mainLayout);
  window->show();
\end{lstlisting}

There are also these three signal functions that they are used to
propagate the click of their respective buttons in other controllers.

\begin{lstlisting}[language=c++, gobble=2]
  void chkFlipXchanged(int arg1);
  void chkFlipYchanged(int arg1);
  void chkSwapXYchanged(int arg1);
\end{lstlisting}

\lvlv{DistancePlotterWidget}

The \textbf{DistancePlotterWidget} class is a widget for plot the
distance between a centroid and a point. This widget is used in
combination with \textbf{CalibratedWiget} and \textbf{ManualWidget}. The
include code is the following.

\begin{lstlisting}[language=c++, gobble=2]
  #include <QApplication>
  #include <PAT/Controller.h>
  #include <PAT/Centroid.h>
  #include <PAT/Utils/Point.h>
  #include <memory>
  #include <thread>
  #include <stdlib.h>
  #include <QDebug>
  using namespace PAT::Controller;
\end{lstlisting}

This is a simple \textbf{main} example.

\begin{lstlisting}[language=c++, gobble=2]
  int main(int argc, char *argv[]) {
    QApplication a(argc, argv);
    DistancePlotterWidget distancePlotterWidget;
    std::shared_ptr<PAT::Centroid> centroid(new PAT::Centroid());
    std::shared_ptr<PAT::Utils::BoundedParameter<PAT::Utils::Circle<int>>> target(new PAT::Utils::BoundedParameter<PAT::Utils::Circle<int>>({0, 1000, 1, 500}, {0, 1000, 1, 500}));
    double frameRate = 30;
    distancePlotterWidget.setCentroidAndTarget(centroid, target);
    std::thread t1([&]() {
      for (size_t i = 0; i < 50000; i++) {
        centroid->x = (double)rand() / RAND_MAX * 1000;
        centroid->y = (double)rand() / RAND_MAX * 1000;
        emit centroid->newValue();
        std::this_thread::sleep_for(std::chrono::milliseconds((int)round(1000 / frameRate)));
      }
    });
    distancePlotterWidget.show();
    t1.detach();
    return a.exec();
  }
\end{lstlisting}

\lvliv{UML}

\img{Controller}{Controller UML.}
            \lvliii{Setup}

This project is a list of apps that they combine custom GUI for the PAT
system.

\textbf{Requirements}

\includegraphics[scale=0.7]{img/shilds/cpp.png}
\includegraphics[scale=0.7]{img/shilds/cmake.png}
\includegraphics[scale=0.7]{img/shilds/git.png}
\includegraphics[scale=0.7]{img/shilds/doxygen.png}
\includegraphics[scale=0.7]{img/shilds/sphinx.png}
\includegraphics[scale=0.7]{img/shilds/win.png}

\lvliv{Generality}

\lvlv{Import}

Import as an external library into your project by copy-paste the
following lines in your \texttt{config.json}.

\begin{lstlisting}[language=javascript, gobble=2]
  {
    "name"     : "PATSetup",
    "path"     : "gitlab.dei.unipd.it/PAT/Setup.git",
    "tag"      : "HEAD",
    "available": "YES",
    "getGui"   : "NO"
  }
\end{lstlisting}

\lvlv{Prerequisites}
The Prerequisites are the same descibed in the section "Setup your enviroment Prerequisites" \sect{Setup your environment} with the addition of the following libraries that are auto fetched from the gitlab.dei.unipd.it:

\begin{itemize}
  \tightlist
  \item
        \href{https://gitlab.dei.unipd.it/PAT/Controller.git}{Controller}
        HEAD
\end{itemize}


\lvliv{Usage}
In this repo you can create your own GUI for the PAT system in the
\texttt{apps} folder. As reference you can use one of the already
defined apps. There is no need to explain additional things because in
this repository you want to compose ready-made pieces and not create new
ones. To simplify your orientation, you can find a global image of all the repositories in the picture \fig{Setup}.

\img{Setup}{Setup UML without attributes.}

      }